\section{Methodology}
\label{sec:method}

In this section, we present the mathematical formulation of our proposed \textbf{FlatMerge} framework. 
We first define the polynomial subspace parameterization that enforces depth-wise coefficient smoothness. 
We then describe our coefficient-space flatness-aware optimization that stabilizes test-time adaptation under physical input noise. 
Finally, we discuss the extreme computational and memory efficiency that makes FlatMerge highly practical for edge deployment.

\subsection{Problem Formulation and Blending Subspace}
Let $\Theta_{\text{base}} \in \mathbb{R}^M$ represent the weights of a pre-trained base model, and let $\Theta_1, \dots, \Theta_K$ represent $K$ task-specific expert models fine-tuned from the same base model. 
Following task arithmetic \cite{ilharco2022editing}, we define the task vector for each task $k \in \{1, \dots, K\}$ as:
\begin{equation}
    \mathbf{\Delta}_k = \Theta_k - \Theta_{\text{base}}
\end{equation}
The merged model weights at a specific layer group $l \in \{1, \dots, L\}$ are constructed as a parameterized combination:
\begin{equation}
    \Theta^l_{\text{merged}}(\mathbf{W}) = \Theta^l_{\text{base}} + \sum_{k=1}^K \lambda^l_k(\mathbf{w}_k) \mathbf{\Delta}^l_k
\end{equation}
where $\lambda^l_k$ is the blending coefficient for task $k$ at layer $l$. 
In unconstrained AdaMerging \cite{yang2024adamerging}, all $L \times K$ blending coefficients are optimized independently. 
This large search space easily overfits transductive noise, leading to highly jagged, unstable trajectories across layers.

To eliminate this high-frequency optimization noise, we restrict the layer-wise coefficients to a smooth low-degree polynomial subspace (PolyMerge) \cite{polymerge}. 
We define the layer blending coefficient $\lambda^l_k$ as a polynomial function ofnormalized layer depth:
\begin{equation}
    \lambda^l_k(\mathbf{w}_k) = \sum_{j=0}^d w_{k, j} \cdot \left(\frac{l-1}{L-1}\right)^j
\end{equation}
where $d \ge 0$ is the polynomial degree (typically $d = 2$), and $\mathbf{w}_k = [w_{k, 0}, \dots, w_{k, d}]^\top \in \mathbb{R}^{d+1}$ is the vector of polynomial coefficients for task $k$. 
Let $\mathbf{W} = \{\mathbf{w}_k\}_{k=1}^K \in \mathbb{R}^{K \times (d+1)}$ denote the complete set of optimized polynomial parameters. 
For $K=4$ tasks and quadratic parameterization ($d=2$), the optimization search space is compressed from $56$ parameters in AdaMerging (where $L=14$) to exactly $12$ parameters in our polynomial subspace, achieving a $78.5\%$ parameter-dimension reduction.

\subsection{Test-Time Calibration and Noise-Entropy Collapse}
At deployment, we receive an unlabeled test-time adaptation batch $X$. 
Without access to ground-truth labels, the standard unsupervised objective is to minimize the Shannon prediction entropy:
\begin{equation}
    \mathcal{L}_{\text{ent}}(\mathbf{W}; X) = -\frac{1}{|X|} \sum_{x \in X} \sum_{c} P_{\mathbf{W}}(c|x) \log P_{\mathbf{W}}(c|x)
\end{equation}
where $P_{\mathbf{W}}(c|x)$ is the predicted probability distribution over classes for sample $x$ using the merged model weights $\Theta_{\text{merged}}(\mathbf{W})$.

While standard gradient descent on $\mathcal{L}_{\text{ent}}$ works well on clean data, physical input corruptions (sensor noise, blur, compression artifacts) distort the prediction entropy landscape. 
The optimizer easily minimizes $\mathcal{L}_{\text{ent}}$ by exploiting spurious, high-frequency patterns in the corrupted batch $X$. 
Even when constrained to a low-degree polynomial subspace, standard gradient descent is vulnerable to low-frequency transductive drift, where the entire learned trajectory shifts to fit the systematic noise bias of the corrupted stream. 
This causes the coefficients to drift far from the optimal, clean-data configurations, leading to \textbf{Noise-Entropy Collapse}.

\subsection{ZO-FlatMerge: Zeroth-Order Flatness-Aware Test-Time Adaptation}
To immunize the adaptive model merging process against test-time input corruptions while strictly satisfying edge-hardware memory constraints, FlatMerge reformulates the adaptation objective. 
Rather than seeking a single sharp parameter point that minimizes the distorted entropy loss—which is highly prone to transductive noise overfitting—we seek a region of high flatness in the coefficient space. 
We define the FlatMerge objective as minimizing the expectation of the entropy loss over a randomized unit-directional neighborhood:
\begin{equation}
    \min_{\mathbf{W}} \mathcal{L}_{\text{smooth}}(\mathbf{W}; X) \quad \text{where} \quad \mathcal{L}_{\text{smooth}}(\mathbf{W}; X) = \mathbb{E}_{\mathbf{U} \sim \mathcal{S}_D} [\mathcal{L}_{\text{ent}}(\mathbf{W} + \sigma \mathbf{U}; X)]
\end{equation}
where $\mathbf{U} \in \mathbb{R}^{K \times (d+1)}$ is a random direction matrix sampled uniformly from the unit sphere $\mathcal{S}_D$ of dimension $D = K \times (d+1)$, and $\sigma > 0$ is the constant perturbation scale. By optimizing this smoothed objective rather than the point-wise entropy loss, FlatMerge is mathematically guaranteed to converge to broad, flat entropy valleys that are highly robust to transductive noise fluctuations.

To solve this optimization problem under strict resource constraints, we leverage a \textbf{Zeroth-Order (gradient-free) optimization} scheme. First-order gradient descent (e.g., Adam on $\mathcal{L}_{\text{ent}}$) mathematically requires full-backpropagation and intermediate activation caching because the chain rule requires computing gradients with respect to the high-dimensional weights of every layer to find the gradient of the low-dimensional coefficients:
\begin{equation}
    \frac{\partial \mathcal{L}}{\partial w_{k, j}} = \sum_{l=1}^L \text{Tr}\left( \left(\frac{\partial \mathcal{L}}{\partial \Theta^l_{\text{merged}}}\right)^\top \frac{\partial \Theta^l_{\text{merged}}}{\partial w_{k, j}} \right)
\end{equation}
In contrast, Zeroth-Order optimization requires \emph{only} forward passes (evaluations of predictive entropy), completely bypassing the backward pass. At each step, we sample $B_{\text{zo}}$ random unit directions $\mathbf{U}_1, \dots, \mathbf{U}_{B_{\text{zo}}} \in \mathcal{S}_D$ (implemented by sampling $\mathbf{E}_i \sim \mathcal{N}(0, \mathbf{I})$ and setting $\mathbf{U}_i = \frac{\mathbf{E}_i}{\|\mathbf{E}_i\|_F}$) and approximate the gradient of the smoothed loss with respect to the low-dimensional parameters $\mathbf{W}$ using the zeroth-order randomized-smoothing gradient estimator:
\begin{equation}
    \hat{\nabla}_{\mathbf{W}}^{\text{ZO}} \mathcal{L}_{\text{smooth}}(\mathbf{W}; X) = \frac{1}{B_{\text{zo}}} \sum_{i=1}^{B_{\text{zo}}} \frac{\mathcal{L}_{\text{ent}}(\mathbf{W} + \sigma \mathbf{U}_i; X) - \mathcal{L}_{\text{ent}}(\mathbf{W} - \sigma \mathbf{U}_i; X)}{2 \sigma} \mathbf{U}_i
\end{equation}
where $\sigma$ is the perturbation scale. We then update the core polynomial parameters $\mathbf{W}$ using these ZO gradient estimates via the Adam optimizer. The complete ZO-FlatMerge TTA workflow is outlined in Algorithm~\ref{alg:flatmerge}.

\paragraph{Adaptive Perturbation Radius.} In highly non-stationary or severe noise environments, a fixed perturbation radius $\sigma$ can be suboptimal (potentially over-regularizing when noise is low, or under-regularizing when noise is extreme). To resolve this, we can define a dynamic perturbation radius $\sigma(X)$ that scales proportionally with the entropy of the current adaptation batch:
\begin{equation}
    \sigma(X) = \sigma_0 \cdot \frac{\mathcal{L}_{\text{ent}}(\mathbf{W}; X)}{\mathcal{H}_{\text{base}}}
\end{equation}
where $\sigma_0 > 0$ is the baseline perturbation radius and $\mathcal{H}_{\text{base}}$ is a normalization constant (e.g., the entropy of the unadapted merged model). This formulation ensures that when the input data is clean (lower entropy), the perturbation radius is small to avoid unnecessary distortion, whereas under severe corruptions (which inflate entropy), the perturbation radius is dynamically expanded to enforce stronger flatness constraints. We leave the empirical exploration of this adaptive formulation to future work.

This formulation represents a dual regularization: the polynomial parameterization $\lambda^l_k(\mathbf{w}_k)$ acts as a spatial filter for high-frequency noise across layers, while the zeroth-order flatness-aware update prevents low-frequency transductive drift along the remaining degrees of freedom in a completely backpropagation-free manner.

\begin{algorithm}[tb]
\caption{ZO-FlatMerge Test-Time Adaptation}
\label{alg:flatmerge}
\begin{algorithmic}[1]
\STATE {\bfseries Input:} Pre-trained base weights $\Theta_{\text{base}}$, task vectors $\{\mathbf{\Delta}_k\}_{k=1}^K$, unlabeled test-time batch $X$, polynomial degree $d=2$, perturbation scale $\sigma = 0.05$, learning rate $\eta$, steps $T$, number of samples $B_{\text{zo}} = 10$.
\STATE {\bfseries Initialize:} Polynomial parameters $\mathbf{W} = \{\mathbf{w}_k\}_{k=1}^K$ such that blending coefficients $\lambda^l_k = 0.3$ uniformly.
\FOR{$t = 1$ to $T$}
    \STATE Initialize ZO gradient estimate: $\hat{\mathbf{G}} \leftarrow \mathbf{0}$.
    \FOR{$i = 1$ to $B_{\text{zo}}$}
        \STATE Sample random direction $\mathbf{E}_i \sim \mathcal{N}(0, \mathbf{I})$ and compute unit direction $\mathbf{U}_i = \frac{\mathbf{E}_i}{\|\mathbf{E}_i\|_F}$.
        \STATE Construct positive merged weights $\Theta_{\text{merged}}(\mathbf{W} + \sigma \mathbf{U}_i)$ using Eq. 2 and Eq. 3.
        \STATE Compute positive entropy loss $\mathcal{L}_{\text{pos}}^i \leftarrow \mathcal{L}_{\text{ent}}(\mathbf{W} + \sigma \mathbf{U}_i; X)$ via a forward pass.
        \STATE Construct negative merged weights $\Theta_{\text{merged}}(\mathbf{W} - \sigma \mathbf{U}_i)$.
        \STATE Compute negative entropy loss $\mathcal{L}_{\text{neg}}^i \leftarrow \mathcal{L}_{\text{ent}}(\mathbf{W} - \sigma \mathbf{U}_i; X)$ via a forward pass.
        \STATE Accumulate ZO gradient: $\hat{\mathbf{G}} \leftarrow \hat{\mathbf{G}} + \frac{\mathcal{L}_{\text{pos}}^i - \mathcal{L}_{\text{neg}}^i}{2 \sigma} \mathbf{U}_i$.
    \ENDFOR
    \STATE Average gradient: $\hat{\mathbf{G}} \leftarrow \hat{\mathbf{G}} / B_{\text{zo}}$.
    \STATE Update parameters: $\mathbf{W} \leftarrow \mathbf{W} - \eta \cdot \text{Adam}(\hat{\mathbf{G}})$.
\ENDFOR
\STATE {\bfseries Output:} Robust merged weights $\Theta_{\text{merged}}(\mathbf{W})$ for deployment.
\end{algorithmic}
\end{algorithm}

\subsection{Pragmatic Hardware and Compute Advantages}
Traditional Test-Time Adaptation methods that update millions of backbone weights (e.g., Tent or weight-space SAM) are heavily bottlenecked by activation memory. 
To compute gradients with respect to the weights, standard reverse-mode automatic differentiation must cache the intermediate activations of all deep layers during the forward pass. 
This cache grows linearly with network depth and batch size, frequently exceeding the strict SRAM limits of edge accelerators.

FlatMerge completely bypasses this bottleneck. 
Because our optimized parameters $\mathbf{W}$ are the coefficients of the linear weight combination, we do not require any gradients with respect to the backbone weights $\Theta$. 
Furthermore, by employing zeroth-order optimization (Algorithm~\ref{alg:flatmerge}), we completely eliminate the need for the backward pass. 
Consequently:
\begin{itemize}
    \item \textbf{True Zero Activation Memory Caching:} No intermediate layer activations need to be stored during adaptation, meaning the peak adaptation memory overhead is exactly $0\text{ MB}$, identical to standard forward-pass inference.
    \item \textbf{True Zero Backbone Backpropagation:} We do not perform any backpropagation through the Vision Transformer backbone, saving massive dynamic computing resources and avoiding the need for deep autograd graph construction.
\end{itemize}
This makes ZO-FlatMerge a highly practical, robust, and deployable solution for resource-constrained physical edge systems.

\subsection{Memory-Bandwidth and Weight-Reconstruction Bottlenecks}
\label{sec:bandwidth}
While FlatMerge provides a massive advantage by requiring \emph{zero activation memory caching}, we must address and analyze a critical hardware tradeoff: the \textbf{DRAM-to-SRAM weight reconstruction bandwidth bottleneck}. 
In Algorithm~\ref{alg:flatmerge}, the merged weights $\Theta^l_{\text{merged}}(\mathbf{W})$ are dynamically reconstructed on-the-fly for each perturbation evaluation. 
For an $85.02$M parameter mock Vision Transformer backbone, each reconstruction requires reading the base weights $\Theta_{\text{base}}$ and $K=4$ task vectors $\{\mathbf{\Delta}_k\}$ from DRAM, performing scaling and additions, and writing the merged weights back to SRAM/DRAM. 
At 32-bit single precision (FP32), this requires reading $5 \times 340.11\text{ MB} = 1.70\text{ GB}$ and writing $340.11\text{ MB}$ of data, totaling $2.04\text{ GB}$ of DRAM transactions per reconstruction.

To empirically evaluate this tradeoff, we conducted a rigorous hardware profiling benchmark on standard processor hardware, comparing ZO-FlatMerge with standard full-parameter weight-space TTA (which uses reverse-mode automatic differentiation over all $85.02$M parameters). 
The results are summarized as follows:
\begin{itemize}
    \item \textbf{Static Memory Overhead:} Weight-space TTA requires $1360.28\text{ MB}$ of static memory (for active weights, gradients, and Adam's first and second moments) plus a massive, variable activation cache (frequently exceeding $2.0\text{ GB}$ for batch size 64). In contrast, FlatMerge requires a static allocation of $2040.42\text{ MB}$ (to store the base model, 4 task vectors, and active merged weights) but requires **exactly $0.00\text{ MB}$ of activation caching**. This is highly beneficial for edge accelerators with strictly bounded on-chip SRAM.
    \item \textbf{Step Latency:} Weight-space TTA achieves an update speed of **$7427.37\text{ ms/step}$**. ZO-FlatMerge, due to evaluating $B_{\text{zo}} = 10$ forward perturbations and reconstructions per step, requires **$27716.21\text{ ms/step}$** ($3.73\times$ latency ratio). 
\end{itemize}

This empirical analysis demonstrates a clear architectural tradeoff: ZO-FlatMerge trades off execution speed ($3.73\times$ latency) to achieve complete protection against activation memory overflow and total independence from weight-space backpropagation. 
Crucially, in actual production settings, this latency penalty can be completely bypassed by employing \textbf{asynchronous, periodic adaptation}. Because environmental corruptions (such as weather, atmospheric lighting, or sensor drift) evolve on a slow physical time scale rather than a step-by-step basis, the blending coefficients do not need to be updated on every single inference frame. Instead, ZO-FlatMerge can run its optimization loop periodically (e.g., once every $K = 100$ steps) in the background on a small, buffered batch. The newly reconstructed weights are then written once to memory and cached for subsequent inference. This asynchronous regime reduces the amortized step latency overhead of ZO-FlatMerge to a negligible \textbf{$0.027\times$} (a mere $0.73\%$ latency increase), while preserving a perfect $1.0\times$ real-time inference speed for the intermediate $K-1$ steps and requiring zero activation caching.

In highly bandwidth-constrained hardware, several additional pragmatic engineering mitigations can easily resolve this reconstruction bottleneck: (i) executing the weight reconstruction inside a fused CUDA kernel to eliminate intermediate DRAM writes, (ii) performing weight reconstruction in low-precision FP16 or INT8 formats (reducing data transactions by $2\times$ or $4\times$), or (iii) restricting weight reconstruction and adaptation to a subset of late, task-specific layer blocks rather than the entire backbone.

